\chapter{Proper actions and the universal proper space}
\label{ch:proper-actions}

This chapter begins with the verified topology of proper discrete-group
actions and their subgroup fixed-point spaces. Continuous equivariant maps,
homotopies, and their fixed-point restrictions are also verified. Equivariant CW structures,
the existence and universal property of $\underline E\Gamma$, and the
directed system of cocompact pieces remain formalization targets.

\section{Proper actions and compact transporters}

\begin{definition}\label{def:proper-action}
\lean{BC4lean.ProperActions.actionMap,BC4lean.ProperActions.proper_iff_actionMap}
\leanok
For a topological group $\Gamma$ acting on a topological space $X$, the
action is proper when
\[
 \alpha:\Gamma\times X\longrightarrow X\times X,
 \qquad \alpha(g,x)=(gx,x)
\]
is a proper map. This is exactly Mathlib's \texttt{ProperSMul} predicate;
it includes continuity of the action. In the applications below $\Gamma$
is discrete. Finite stabilizers alone do not replace properness.
\end{definition}

\begin{definition}\label{def:compact-transporter}
\lean{BC4lean.ProperActions.transporter}
\leanok
For subsets $K,L\subseteq X$, define the transporter
\[
 T(K,L)=\{g\in\Gamma:gK\cap L\ne\varnothing\}.
\]
\end{definition}

\begin{proposition}\label{prop:compact-transporter-criterion}
\uses{def:proper-action,def:compact-transporter}
\lean{BC4lean.ProperActions.finite_transporter,BC4lean.ProperActions.proper_iff_finite_transporters,BC4lean.ProperActions.proper_iff_compact_preimages}
For a proper action of a discrete group, $T(K,L)$ is finite whenever
$K,L$ are compact. If $X$ is locally compact Hausdorff and the action is
continuous, the converse holds. Under these hypotheses properness is also
equivalent to compactness of $\alpha^{-1}(C)$ for every compact
$C\subseteq X\times X$.
\end{proposition}
\begin{proof}
\leanok
The first projection maps the compact set $\alpha^{-1}(L\times K)$
onto $T(K,L)$. Compact subsets of a discrete space are finite.
Conversely, for compact $C\subseteq X\times X$, its coordinate projections
are compact. The inverse image of $C$ is a closed subset of the product
of a finite transporter and the second coordinate projection.
The pinned Mathlib proper-map and compact-transporter criteria formalize
these implications on locally compact Hausdorff spaces.
\end{proof}

\begin{proposition}\label{prop:proper-finite-stabilizers}
\uses{prop:compact-transporter-criterion}
\lean{BC4lean.ProperActions.finite_stabilizer}
Every point stabilizer $\Gamma_x$ in a proper discrete-group action is
finite. This implication does not require local compactness of $X$.
\end{proposition}
\begin{proof}
\leanok
The stabilizer is $T(\{x\},\{x\})$, and singletons are compact.
\end{proof}

\begin{proposition}\label{prop:proper-orbit-space}
\uses{def:proper-action}
\lean{BC4lean.ProperActions.orbitSpace_t2,BC4lean.ProperActions.orbitSpace_locallyCompact}
A proper action has Hausdorff orbit space $X/\Gamma$. If $X$ is locally
compact, then $X/\Gamma$ is locally compact as well.
\end{proposition}
\begin{proof}
\leanok
Use Mathlib's Hausdorff quotient theorem for proper actions. The orbit
projection is an open quotient map for a continuous group action, and
local compactness descends through an open quotient map.
\end{proof}

\begin{proposition}\label{prop:proper-subgroup-restriction}
\uses{def:proper-action}
\lean{BC4lean.ProperActions.proper_subgroup}
Restricting a proper discrete-group action to any subgroup preserves
properness.
\end{proposition}
\begin{proof}
\leanok
Every subgroup of a discrete group is closed. Apply the proper-action
restriction theorem for a closed embedding of groups.
\end{proof}

\begin{proposition}\label{prop:proper-action-examples}
\uses{prop:compact-transporter-criterion}
\lean{BC4lean.ProperActions.proper_of_finite,BC4lean.ProperActions.proper_leftTranslation,BC4lean.ProperActions.finite_group_of_compact}
Every continuous action of a finite discrete group on a locally compact
Hausdorff space is proper. Left translation of a topological group on
itself is proper. If a discrete group acts properly on a nonempty compact
space, the group is finite.
\end{proposition}
\begin{proof}
\leanok
For a finite group every transporter is finite. For left translation,
$(g,x)\mapsto(gx,x)$ is a homeomorphism with inverse
$(y,x)\mapsto(yx^{-1},x)$. For the last assertion,
$T(X,X)=\Gamma$ is finite by compactness and nonemptiness.
\end{proof}

\section{Subgroup fixed-point spaces}

\begin{definition}\label{def:subgroup-fixed-point-space}
\lean{BC4lean.ProperActions.FixedPointSpace}
\leanok
For a subgroup $H\leq\Gamma$, let
\[
 X^H=\{x\in X: hx=x\text{ for every }h\in H\}
\]
with the subspace topology. The Lean definition uses Mathlib's existing
\texttt{MulAction.fixedPoints}.
\end{definition}

\begin{proposition}\label{prop:fixed-point-subgroup-laws}
\uses{def:subgroup-fixed-point-space}
\lean{BC4lean.ProperActions.mem_fixedPoints_iff_le_stabilizer,BC4lean.ProperActions.fixedPoints_antitone,BC4lean.ProperActions.fixedPoints_bot}
A point $x$ lies in $X^H$ exactly when $H\leq\Gamma_x$.
If $H\leq K$, then $X^K\subseteq X^H$, and $X^{\{1\}}=X$.
\end{proposition}
\begin{proof}
\leanok
Unfold the fixed-point and stabilizer conditions and use subgroup inclusion.
\end{proof}

\begin{proposition}\label{prop:fixed-point-space-closed}
\uses{def:subgroup-fixed-point-space}
\lean{BC4lean.ProperActions.isClosed_fixedPoints,BC4lean.ProperActions.fixedPointSpace_locallyCompact}
For an action by continuous translations on a Hausdorff space, $X^H$
is closed for every subgroup $H$. If $X$ is locally compact, so is $X^H$.
\end{proposition}
\begin{proof}
\leanok
The fixed-point set is the intersection, over $h\in H$, of the closed
equalizers of $x\mapsto hx$ and $x\mapsto x$. Closed subspaces of locally
compact spaces are locally compact.
\end{proof}

\begin{proposition}\label{prop:proper-infinite-fixed-points}
\uses{prop:proper-finite-stabilizers,prop:fixed-point-subgroup-laws}
\lean{BC4lean.ProperActions.finite_subgroup_of_fixedPoint,BC4lean.ProperActions.fixedPoints_eq_empty_of_infinite,BC4lean.ProperActions.fixedPointSpace_isEmpty}
For a proper discrete-group action, any subgroup fixing a point is finite.
Consequently $X^H=\varnothing$ whenever $H$ is infinite.
\end{proposition}
\begin{proof}
\leanok
If $x\in X^H$, then $H$ is a subgroup of the finite group $\Gamma_x$.
\end{proof}

Properness does not imply contractibility of fixed-point spaces for finite
subgroups. That is additional structure required of a universal proper model.

\section{Equivariant maps and homotopies}

\begin{definition}\label{def:equivariant-map}
\lean{BC4lean.ProperActions.EquivariantMap,BC4lean.ProperActions.EquivariantMap.toMulActionHom}
\leanok
For topological spaces $X,Y$ with $\Gamma$-actions, a continuous equivariant
map $f:X\to Y$ is a continuous map satisfying $f(gx)=gf(x)$.
The underlying equivariant function uses Mathlib's \texttt{MulActionHom}.
The constructions in this section do not require properness, a discrete
group topology, local compactness, or Hausdorffness.
\end{definition}

\begin{proposition}\label{prop:equivariant-map-laws}
\uses{def:equivariant-map}
\lean{BC4lean.ProperActions.EquivariantMap.ext,BC4lean.ProperActions.EquivariantMap.map_smul,BC4lean.ProperActions.EquivariantMap.continuous,BC4lean.ProperActions.EquivariantMap.id,BC4lean.ProperActions.EquivariantMap.comp,BC4lean.ProperActions.EquivariantMap.id_apply,BC4lean.ProperActions.EquivariantMap.comp_apply,BC4lean.ProperActions.EquivariantMap.id_comp,BC4lean.ProperActions.EquivariantMap.comp_id,BC4lean.ProperActions.EquivariantMap.comp_assoc}
Continuous equivariant maps have identity and associative composition, with
both identity laws. Equality of maps is determined pointwise.
\end{proposition}
\begin{proof}
\leanok
Use the identity and composition of continuous maps. Equivariance is
preserved because $g(f(ax))=g(af(x))=ag(f(x))$. The laws follow pointwise.
\end{proof}

\begin{proposition}\label{prop:fixed-point-map-functoriality}
\uses{prop:equivariant-map-laws,def:subgroup-fixed-point-space}
\lean{BC4lean.ProperActions.EquivariantMap.map_fixedPoint,BC4lean.ProperActions.EquivariantMap.fixedPointsMap,BC4lean.ProperActions.EquivariantMap.fixedPointsMap_apply,BC4lean.ProperActions.EquivariantMap.fixedPointsMap_id,BC4lean.ProperActions.EquivariantMap.fixedPointsMap_comp}
For every subgroup $H\leq\Gamma$, an equivariant map $f:X\to Y$ restricts
to a continuous map $f^H:X^H\to Y^H$. Restriction preserves identity and
composition. The subgroup need not be normal.
\end{proposition}
\begin{proof}
\leanok
For $h\in H$ and $x\in X^H$, one has $hf(x)=f(hx)=f(x)$.
Continuity follows from the subspace topology; the laws follow pointwise.
\end{proof}

\begin{definition}\label{def:equivariant-homotopy}
\uses{def:equivariant-map}
\lean{BC4lean.ProperActions.EquivariantHomotopy,BC4lean.ProperActions.EquivariantlyHomotopic}
\leanok
An equivariant homotopy from $f_0$ to $f_1$ is a jointly continuous map
$F:[0,1]\times X\to Y$ with $F(0,x)=f_0(x)$, $F(1,x)=f_1(x)$, and
$F(t,gx)=gF(t,x)$. The interval has trivial action.
Write $f_0\simeq_\Gamma f_1$ when such a homotopy exists.
The Lean definition specializes Mathlib's \texttt{ContinuousMap.HomotopyWith}.
\end{definition}

\begin{proposition}\label{prop:equivariant-homotopy-operations}
\uses{def:equivariant-homotopy,prop:equivariant-map-laws}
\lean{BC4lean.ProperActions.EquivariantHomotopy.refl,BC4lean.ProperActions.EquivariantHomotopy.symm,BC4lean.ProperActions.EquivariantHomotopy.trans,BC4lean.ProperActions.EquivariantHomotopy.comp,BC4lean.ProperActions.EquivariantHomotopy.atTime,BC4lean.ProperActions.EquivariantHomotopy.at_apply,BC4lean.ProperActions.EquivariantHomotopy.at_zero,BC4lean.ProperActions.EquivariantHomotopy.at_one,BC4lean.ProperActions.EquivariantlyHomotopic.refl,BC4lean.ProperActions.EquivariantlyHomotopic.symm,BC4lean.ProperActions.EquivariantlyHomotopic.trans,BC4lean.ProperActions.EquivariantlyHomotopic.comp}
Equivariant homotopy is reflexive, symmetric, transitive, and compatible
with composition. Every time slice of an equivariant homotopy is a
continuous equivariant map, with the specified maps at times zero and one.
\end{proposition}
\begin{proof}
\leanok
Use constant homotopies, time reversal, and concatenation on the two half
intervals. Composition is $(G\circ F)(t,x)=G(t,F(t,x))$.
Joint continuity follows from Mathlib's homotopy operations; equivariance
holds at each time. Restriction to a time slice is continuous.
\end{proof}

\begin{proposition}\label{prop:fixed-point-homotopy-restriction}
\uses{prop:equivariant-homotopy-operations,prop:fixed-point-map-functoriality}
\lean{BC4lean.ProperActions.EquivariantHomotopy.onFixedPoints,BC4lean.ProperActions.EquivariantHomotopy.onFixedPoints_apply,BC4lean.ProperActions.EquivariantHomotopy.onFixedPoints_refl,BC4lean.ProperActions.EquivariantHomotopy.onFixedPoints_symm,BC4lean.ProperActions.EquivariantHomotopy.onFixedPoints_trans,BC4lean.ProperActions.EquivariantHomotopy.onFixedPoints_comp,BC4lean.ProperActions.EquivariantlyHomotopic.onFixedPoints}
An equivariant homotopy $F:f_0\simeq_\Gamma f_1$ restricts to an ordinary
continuous homotopy $F^H:f_0^H\simeq f_1^H$ for every subgroup $H$.
Restriction preserves constant homotopies, reversal, concatenation with
the chosen half-interval parametrization, and composition of homotopies.
\end{proposition}
\begin{proof}
\leanok
At each time equivariance ensures $F(t,X^H)\subseteq Y^H$.
The map into $Y^H$ is jointly continuous by the subspace topology, and the
endpoint identities hold as equalities of subtype elements. The operation
laws follow from their formulas, splitting the two time intervals for
concatenation.
\end{proof}

\begin{definition}\label{def:equivariant-homotopy-equivalence}
\uses{def:equivariant-homotopy,prop:equivariant-map-laws}
\lean{BC4lean.ProperActions.EquivariantHomotopyEquiv}
\leanok
An equivariant homotopy equivalence between $X$ and $Y$ consists of
continuous equivariant maps $f:X\to Y$ and $g:Y\to X$ with
$gf\simeq_\Gamma\operatorname{id}_X$ and
$fg\simeq_\Gamma\operatorname{id}_Y$.
\end{definition}

\begin{proposition}\label{prop:fixed-point-homotopy-equivalence}
\uses{def:equivariant-homotopy-equivalence,prop:fixed-point-homotopy-restriction}
\lean{BC4lean.ProperActions.EquivariantHomotopyEquiv.refl,BC4lean.ProperActions.EquivariantHomotopyEquiv.symm,BC4lean.ProperActions.EquivariantHomotopyEquiv.trans,BC4lean.ProperActions.EquivariantHomotopyEquiv.fixedPoints,BC4lean.ProperActions.EquivariantHomotopyEquiv.fixedPoints_toFun_apply,BC4lean.ProperActions.EquivariantHomotopyEquiv.fixedPoints_invFun_apply}
Equivariant homotopy equivalences have identity, reversal, and composition.
Every equivariant homotopy equivalence $X\simeq_\Gamma Y$ induces an
ordinary homotopy equivalence $X^H\simeq Y^H$ for every subgroup $H$.
\end{proposition}
\begin{proof}
\leanok
Compose the equivariant maps and their inverse homotopies to obtain the
three operations. Restrict both maps and both inverse homotopies to the
fixed-point spaces, using the proved identity and composition laws.
\end{proof}

\section{Equivariant CW complexes and a model for $\underline E\Gamma$}

\begin{definition}\label{def:equivariant-cw-complex}
\uses{def:proper-action,def:equivariant-map,def:equivariant-homotopy}
\notready
Using the preceding maps and homotopies, develop $\Gamma$-CW complexes by
equivariant cell attachment, with cells $\Gamma/H\times D^n$.
For the proper theory the cell stabilizers $H$ are finite. Establish the
connection between this cellular condition and topological properness.
\end{definition}

\begin{definition}\label{def:universal-proper-space}
\uses{def:equivariant-cw-complex,def:subgroup-fixed-point-space,prop:proper-infinite-fixed-points}
\notready
A universal proper $\Gamma$-CW complex has contractible fixed-point space
for every finite subgroup and empty fixed-point space for every infinite
subgroup. Develop existence, the universal property, and uniqueness up to
equivariant homotopy. Only the empty-fixed-point consequence of properness
has been verified above; no universal proper model is constructed here.
\end{definition}

\section{Cocompact pieces and model choices}
\begin{definition}\label{def:cocompact-pieces}
\uses{def:universal-proper-space}
\notready
Choose a suitable locally compact second-countable model $E$ for
$\underline E\Gamma$. Index its closed invariant subsets $X$ with compact
quotient $X/\Gamma$. Justify the model choice and the directed indexing system.
\end{definition}
